\section{Integrasi Flex--Bison dan Build}
\label{sec:ch13-integrasi}

\subsection{Antarmuka \code{yylex}}

Bison memanggil \code{yylex} untuk mendapatkan token berikutnya.
Fungsi ini harus mengisi \code{yylval} dan mengembalikan kode token
Bison. Karena \code{lexer\_next} mengembalikan \code{Token}, driver
perlu memetakan jenis token ke kode Bison.

\begin{figure}[H]
\centering
\begin{tikzpicture}[node distance=6mm and 9mm]
  \node[phase, text width=2.6cm] (src) {Program\\sumber};
  \node[phase, text width=2.6cm, right=of src] (flex) {Flex\\\code{yylex}};
  \node[phase, text width=2.6cm, right=of flex] (bison)
        {Bison\\\code{yyparse}};
  \node[phase, text width=2.6cm, below=of bison] (ast) {AST};
  \draw[arrow] (src) -- (flex);
  \draw[arrow] (flex) -- node[above, font=\small] {token} (bison);
  \draw[arrow] (bison) -- (ast);
\end{tikzpicture}
\caption{Integrasi Flex dan Bison}
\label{fig:ch13-integrasi}
\end{figure}

\subsection{Langkah Build}

Build lengkap:
\begin{lstlisting}[language=bash]
flex -o build/lexer_flex.c src/lexer_flex.l
bison -d -o build/parser.tab.c src/parser.y
gcc -Iinclude -Ibuild build/lexer_flex.c build/parser.tab.c \
    src/token.c src/error.c src/ast.c src/main_bison.c \
    -o build/scc
\end{lstlisting}

Perintah pertama menghasilkan lexer C. Perintah kedua menghasilkan
parser dan header token. Perintah ketiga mengompilasi semuanya
menjadi executable \code{scc}.

\subsection{Driver}

Driver membuka berkas, menyiapkan lexer global, lalu memanggil
\code{yyparse}. Jika parsing berhasil dan tidak ada kesalahan
leksikal, driver mencetak AST atau melanjutkan ke analisis semantik.

\begin{lstlisting}[language=C]
int main(int argc, char **argv) {
    if (argc != 2) {
        fprintf(stderr, "penggunaan: %s berkas.sc\n", argv[0]);
        return 2;
    }
    if (!lexer_open(argv[1])) {
        return 2;
    }
    int status = yyparse();
    lexer_close();
    return status;
}
\end{lstlisting}

\subsection{Catatan Lingkungan}

Lingkungan penyusunan buku ini tidak memiliki Flex, Bison, dan GCC.
Kode pada bab ini disusun sebagai artefak yang dapat dibangun pada
lingkungan dengan toolchain lengkap. Mahasiswa harus menjalankan
build dan tes sendiri, lalu mencatat versi alat yang dipakai.

\begin{warningbox}
Perbedaan versi Bison dapat mengubah jumlah konflik dan nama simbol
yang dihasilkan. Selalu catat versi Bison pada laporan praktikum.
\end{warningbox}
